\documentclass{article}
\usepackage{times}
\usepackage{amsmath}
\usepackage{graphicx}
\usepackage{booktabs}

\title{Genetic Algorithm-Based Fairness Violation Discovery}
\author{Student Name}
\date{}

\begin{document}
\maketitle

\section{Introduction}
Machine learning models used in decision support can exhibit unfair behavior across sensitive attributes.
In this project, we target fairness testing for tabular binary classifiers and investigate whether a search-based
approach can discover discriminatory cases more effectively than random search.
The selected problem is \textit{AI Model Fairness Testing} from the module baseline list.

\section{Related Work}
Fairness testing and bias discovery are typically approached in two ways:
\textit{(i)} group-level fairness metrics (e.g., demographic parity, equalized odds) and
\textit{(ii)} individual-level unfairness discovery by perturbing sensitive attributes.
Search-based software engineering methods (e.g., random search, evolutionary search) are suitable when
the input space is large and objective functions are non-linear.
Random search is simple and easy to reproduce, but often inefficient under limited testing budget.
Genetic algorithms can exploit high-fitness areas via selection and crossover, but may incur higher runtime and
require careful parameter tuning.

\section{Problem Definition}
Let $f(x)$ denote a classifier and $T_s(x)$ the sensitive attribute flip operator.
\begin{equation}
D(x) = |P(y=1|x) - P(y=1|T_s(x))|
\end{equation}
where $D(x)$ is the individual discrimination score. A higher value indicates stronger potential unfairness.

\section{Method}
\subsection{Genetic Algorithm}
Population:
\[
\mathcal{P}^{(t)} = \{x_1, ..., x_N\}
\]
Objective:
\[
\max_x D(x)
\]
The implementation uses tournament selection, uniform crossover, and mutation.
For comparability, we enforce an \textbf{equal evaluation budget} between GA and Random Search.
If population size is $N$ and generations are $G$, then each GA trial evaluates $N \times G$ candidates.
Random Search is configured with the same number of evaluations per trial.

\section{Experiments}
\subsection{Setup}
To improve validity and robustness while keeping the same problem definition and baseline, we evaluate
the tool under the following controlled setup:
\begin{itemize}
    \item Datasets: Adult, COMPAS, German Credit.
    \item Sensitive attributes: \texttt{sex} (Adult), \texttt{race} (COMPAS), \texttt{sex} (German).
    \item Predictive model: RandomForestClassifier.
    \item Number of independent trials: 30 (default).
    \item Budget per trial per method: $B = N \times G$ evaluations (matched budget between GA and Random Search).
    \item Baseline: Random Search.
    \item Fairness metrics:
    \begin{itemize}
        \item Individual discrimination count per trial ($D(x) > \tau$, $\tau=0.2$),
        \item Mean individual discrimination score per trial,
        \item Demographic parity difference,
        \item Equalized odds difference (proxy over generated samples).
    \end{itemize}
\end{itemize}

\subsection{Procedure}
For each dataset, GA and Random Search are executed with identical evaluation budgets and seeds.
For each trial, all fairness metrics are recorded for both methods.
Additionally, convergence traces are captured as:
\begin{itemize}
    \item best fitness per generation/chunk,
    \item mean fitness per generation/chunk.
\end{itemize}
Hyperparameter sensitivity is analyzed by varying:
\begin{itemize}
    \item population size $\in \{10, 20, 40\}$,
    \item mutation rate $\in \{0.05, 0.1, 0.2\}$,
\end{itemize}
while preserving the same total evaluation budget.

\section{Statistical Test}
Wilcoxon signed-rank tests are used at $\alpha = 0.05$ for paired trial results.
Alternative hypothesis: GA $>$ Random Search.
Paired tests are performed for each dataset and each fairness metric:
\begin{itemize}
    \item individual discrimination count,
    \item mean individual discrimination score,
    \item demographic parity difference,
    \item equalized odds difference (proxy).
\end{itemize}

\section{Results}
Quantitative results should be reported from:
\texttt{experiments/results/combined\_results.json} and
\texttt{experiments/stat\_test\_summary.txt}.
Table~\ref{tab:results-template} reports the current run (\textbf{3 trials, budget=1000}).

\begin{table}[h]
    \centering
    \caption{GA vs Random Search across datasets (current run: 3 trials)}
    \label{tab:results-template}
    \begin{tabular}{llcc}
        \toprule
        Dataset & Metric & GA (mean$\pm$std) & Random (mean$\pm$std) \\
        \midrule
        Adult & ID count & $339.3333 \pm 240.3946$ & $0.3333 \pm 0.4714$ \\
        Adult & ID mean & $0.1634 \pm 0.0062$ & $0.0424 \pm 0.0007$ \\
        Adult & DP diff & $0.0142 \pm 0.0148$ & $0.0092 \pm 0.0047$ \\
        Adult & EO diff (proxy) & $0.4598 \pm 0.2411$ & $0.0844 \pm 0.0305$ \\
        \midrule
        COMPAS & ID count & $724.3333 \pm 61.5702$ & $5.3333 \pm 1.8856$ \\
        COMPAS & ID mean & $0.3191 \pm 0.0168$ & $0.0526 \pm 0.0007$ \\
        COMPAS & DP diff & $0.6481 \pm 0.0453$ & $0.0146 \pm 0.0020$ \\
        COMPAS & EO diff (proxy) & $0.7702 \pm 0.1650$ & $0.2697 \pm 0.0674$ \\
        \midrule
        German & ID count & $0.0000 \pm 0.0000$ & $0.0000 \pm 0.0000$ \\
        German & ID mean & $0.0919 \pm 0.0062$ & $0.0268 \pm 0.0003$ \\
        German & DP diff & $0.2266 \pm 0.1583$ & $0.0394 \pm 0.0284$ \\
        German & EO diff (proxy) & $0.6320 \pm 0.1853$ & $0.1483 \pm 0.0595$ \\
        \bottomrule
    \end{tabular}
\end{table}

Wilcoxon signed-rank test results should be summarized per metric and dataset:
\begin{itemize}
    \item Adult: ID count ($p=0.0899$), ID mean ($p=0.125$), DP ($p=0.625$), EO proxy ($p=0.125$).
    \item COMPAS: ID count ($p=0.125$), ID mean ($p=0.125$), DP ($p=0.125$), EO proxy ($p=0.125$).
    \item German: ID count ($p=1.0$, all-zero difference), ID mean ($p=0.125$), DP ($p=0.125$), EO proxy ($p=0.125$).
    \item At $\alpha=0.05$, none of the above are statistically significant in this small-sample run.
\end{itemize}

Convergence and sensitivity analyses should be discussed with generated plots:
\begin{itemize}
    \item \texttt{experiments/fairness\_ga\_results\_<dataset>.png},
    \item sensitivity heatmap embedded in the same figure (when enabled).
\end{itemize}
Despite non-significant tests with $n=3$, GA shows consistently larger mean fairness-gap metrics than Random
in most settings, especially on COMPAS. A larger trial count (e.g., 30) is required for stronger inferential claims.

\section{Reflection}
Current limitations include:
\begin{itemize}
    \item generated-sample based fairness estimation may deviate from real deployment distributions,
    \item equalized odds is computed as a proxy due to unlabeled synthetic candidates,
    \item GA performance depends on hyperparameter settings and search-space encoding.
\end{itemize}
Future improvements include multi-objective optimization, constraint-aware sampling, and
real-data manifold guided candidate generation.

\section{Conclusion}
This project provides a reproducible fairness-testing toolchain that compares GA and Random Search under matched
evaluation budgets, repeated trials, multiple datasets, and multiple fairness metrics.
Final claims about superiority must be made from paired statistical tests and supported by convergence and
sensitivity evidence.

\section{Artifact}
Code repository: \texttt{<REPLACE-WITH-PUBLIC-REPO-URL>}\\
Raw result files are generated under:
\texttt{experiments/results/}.\\
Replication commands:
\begin{itemize}
    \item \texttt{python scripts/download\_adult.py}
    \item \texttt{python scripts/download\_benchmark\_datasets.py}
    \item \texttt{python experiments/run\_experiment.py --with-sensitivity}
    \item \texttt{python experiments/statistical\_test.py}
    \item \texttt{python experiments/visualization.py --dataset adult}
\end{itemize}

\section{References}
[1] S. Bird et al., ``Fairness and Machine Learning,'' fairmlbook.org.
[2] J. H. Holland, \textit{Adaptation in Natural and Artificial Systems}.
[3] F. Wilcoxon, ``Individual comparisons by ranking methods,'' \textit{Biometrics Bulletin}, 1945.

\end{document}